\chapter{DDMRP : principe, dimensionnement et temps}\label{ch:ddmrp}

Le DDMRP (\textit{Demand Driven Material Requirements Planning}) est une méthode de gestion
des matières proposée par C.~Ptak et C.~Smith~\cite{ptak}. Elle se situe dans la lignée du MRP
(\textit{Material Requirements Planning})~\cite{orlicky}, mais elle pose une autre question.
Le MRP calcule des besoins à partir d'un programme de production prévisionnel. Le DDMRP
protège la matière par des tampons (\textit{buffers}) placés à des points stratégiques, et
déclenche les réapprovisionnements à partir de la consommation réelle et de la demande
qualifiée. Il répond à la question : \emph{combien de matière faut-il protéger ou
réapprovisionner maintenant ?} Il ne décide pas quelle commande client doit recevoir cette
matière ; cette décision relève de la logique de commande (\cref{ch:mto-ddmrp}).

\section{Quatre grandeurs}

\begin{description}
  \item[On Hand (stock physique)] la quantité réellement présente en stock.
  \item[Open Supply (approvisionnements ouverts)] la quantité commandée au fournisseur et non
    encore reçue : ce qui est déjà en route.
  \item[Qualified Demand (demande qualifiée)] la demande non encore satisfaite que le buffer
    doit couvrir. Elle exclut la demande future trop lointaine, sauf les pics.
  \item[Net Flow Position (NFP, position de flux net)] la grandeur de décision :
    \begin{equation}\label{eq:nfp}
      NFP = OnHand + OpenSupply - QualifiedDemand.
    \end{equation}
\end{description}

Le stock physique seul ne suffit pas à décider. Avec un stock nul, une commande fournisseur
de 500 unités ouverte et une demande qualifiée de 125 unités, \cref{eq:nfp} donne
$NFP = 0 + 500 - 125 = 375$. La position réelle est donc de 375 unités, dont 125 sont déjà
promises, et une nouvelle commande ne se justifie pas (\cref{fig:nfp}).

\begin{figure}[H]
  \centering
  \begin{tikzpicture}[font=\small]
    \draw[ink, fill=white] (0,0) rectangle (1.6,0.001);
    \node[below] at (0.8,-0.05) {On Hand};
    \node[anchor=west] at (1.8,0.0) {\small $0$};
    \draw[osc, fill=osc!25] (3,0) rectangle (4.6,3.5);
    \node[below] at (3.8,-0.05) {Open Supply};
    \node[anchor=west] at (4.8,1.75) {\small $500$};
    \draw[mtoc, fill=mtoc!30] (6,0) rectangle (7.6,0.875);
    \node[below] at (6.8,-0.05) {Qualified Demand};
    \node[anchor=west] at (7.8,0.44) {\small $125$};
    \node[anchor=west, font=\small\bfseries] at (9.2,1.2) {NFP = 0 + 500 - 125 = 375};
    \node at (2.4,1.8) {\Large $+$};
    \node at (5.3,1.8) {\Large $-$};
  \end{tikzpicture}
  \caption{Position de flux net dans un cas simple : l'approvisionnement ouvert compense
    la demande qualifiée.}
  \label{fig:nfp}
\end{figure}


\section{Buffers rouge, jaune et vert}

Un buffer découpe la plage de NFP en trois zones séparées par deux seuils : le \emph{top of
red} (TOR) et le \emph{top of yellow} (TOY). Au-dessus du TOY, la zone verte indique une
position saine. Entre TOR et TOY, la zone jaune signale qu'un réapprovisionnement doit être
planifié. Sous le TOR, la zone rouge exige une action immédiate (\cref{fig:buffer}).

\begin{figure}[H]
  \centering
  \begin{tikzpicture}[font=\small]
    \fill[ddmrpgrn!35] (0,4.2) rectangle (3,6.0);
    \fill[ddmrpyel!45] (0,2.4) rectangle (3,4.2);
    \fill[ddmrpred!35] (0,0) rectangle (3,2.4);
    \draw[ink, thick] (0,0) rectangle (3,6.0);
    \node[left] at (0,4.2) {\small TOY};
    \node[left] at (0,2.4) {\small TOR};
    \node[left] at (0,6.0) {\small TOG};
    \node at (1.5,5.1) {zone verte};
    \node at (1.5,3.3) {zone jaune};
    \node at (1.5,1.2) {zone rouge};
    \draw[-{Stealth[length=2mm]}, ink] (3.4,0.2) -- (3.4,5.8);
    \node[rotate=90, anchor=south] at (3.9,3.0) {NFP croissant};
    \node[right, text width=5.4cm, align=left] at (4.1,5.1)
      {Position saine : aucune commande.};
    \node[right, text width=5.4cm, align=left] at (4.1,3.3)
      {Planifier un réapprovisionnement jusqu'au TOG.};
    \node[right, text width=5.4cm, align=left] at (4.1,1.2)
      {Action immédiate, priorité maximale.};
  \end{tikzpicture}
  \caption{Buffer DDMRP : zones verte, jaune et rouge, séparées par TOY et TOR.}
  \label{fig:buffer}
\end{figure}


\section{Paramètres et formules}

Le dimensionnement utilise six paramètres : l'ADU (\textit{Average Daily Usage}, consommation
moyenne par heure métier), le lead time LT (en heures métier, durée de réception et de contrôle
incluse), le facteur de lead time LTF, le facteur de variabilité VF, le MOQ (quantité minimale
de commande, multiple auquel on arrondit) et le cycle de commande \emph{orderCycle}.

Les formules suivantes sont celles du noyau implémenté dans SCONTO-SVU. Elles ne doivent pas
être confondues avec la présentation de la méthode dans la littérature.

\begin{align}
  \text{Jaune}          &= ADU \times LT \label{eq:jaune}\\
  \text{Base rouge}     &= ADU \times LT \times LTF \label{eq:base}\\
  \text{Sécurité rouge} &= \text{Base rouge} \times VF \label{eq:secu}\\
  \text{Rouge (TOR)}    &= \text{Base rouge} + \text{Sécurité rouge} \label{eq:rouge}\\
  \text{Verte}          &= \max\bigl(MOQ,\; ADU \times orderCycle,\; ADU \times LT \times LTF\bigr) \label{eq:verte}\\
  \text{TOY}            &= \text{Rouge} + \text{Jaune} \label{eq:toy}\\
  \text{TOG}            &= \text{TOY} + \text{Verte} \label{eq:tog}
\end{align}

La décision suit une règle à deux seuils. Si $NFP \le TOR$, la position est rouge ; si
$TOR < NFP \le TOY$, elle est jaune ; sinon, elle est verte. Lorsque $NFP \le TOY$, une
commande est passée pour $TOG - NFP$, arrondie au multiple supérieur du MOQ.

\begin{table}[H]
  \centering
  \caption{Statut des formules : littérature DDMRP et implémentation SCONTO-SVU.}\label{tab:statut-formules}
  \begin{tabularx}{\textwidth}{@{}YY@{}}
    \toprule
    Élément & Statut \\
    \midrule
    Définition des quatre grandeurs et de la NFP (\cref{eq:nfp}) & concept de la méthode DDMRP~\cite{ptak} \\
    Principe des trois zones et de la décision par seuils & concept de la méthode DDMRP~\cite{ptak} \\
    Formules \cref{eq:jaune} à \cref{eq:tog} & \textbf{implémentation SCONTO-SVU} : valeurs à consolider bibliographiquement \\
    Arrondi au multiple du MOQ & \textbf{implémentation SCONTO-SVU} \\
    Valeurs des paramètres dans les fixtures & \textbf{paramètres de test}, non industriels \\
    \bottomrule
  \end{tabularx}
\end{table}

Une note de méthode s'impose. Le choix du facteur de lead time, du facteur de variabilité et
de la composante \emph{orderCycle} doit encore être consolidé bibliographiquement, puis calibré
avec des données industrielles. Les paramètres des fixtures ne sont pas des paramètres réels
de ZENER.

\section{Temps métier et temps de simulation}

Le modèle a deux horloges : celle du moteur AnyLogic, en secondes simulées, et l'horloge
métier, où un lead time de six heures dure six heures. Le lien est le facteur
\jv{simToRealSeconds} : une seconde simulée représente un nombre donné de secondes réelles.
L'horloge métier utilisée par le DDMRP est
\[
heuresMetier(t) = \frac{(t - tempsDebutSimulation)\times simToRealSeconds}{3600},
\]
où \jv{tempsDebutSimulation} est l'instant de début effectif du run. Le temps écoulé avant ce
début, par exemple pendant le démarrage interactif, n'entre donc pas dans l'ADU.

L'ADU, le lead time et les zones sont calculés en temps métier. La fonction \jv{dureeEchelle()}
sert seulement à planifier les événements dans AnyLogic, par exemple l'arrivée d'une réception.
Modifier \jv{simToRealSeconds} accélère ou ralentit le déroulement du run sans changer le sens
économique du buffer : à état métier identique, les zones ne changent pas (\cref{ch:cas}).
Les délais de promesse client suivent la même règle : ils sont exprimés en secondes métier et
convertis une seule fois en instant de simulation (\jv{echeancePromesse}).
